% Source PDF page 80; manual page 2-51.
\section{Contributions to Acceleration}\label{sec:acceleration-contributions}

The vehicle acceleration, given by the equations of motion, is
\begin{equation}
  \vec a_{\oplus}
  =\frac{\sum\vec F}{m}
   -2\left({}_{I}\vec\omega_{\oplus}\times\vec V_{\oplus}\right)
   -{}_{I}\vec\omega_{\oplus}\times
      \left({}_{I}\vec\omega_{\oplus}\times\vec r\right).
  \label{eq:acceleration-force-contributions}
\end{equation}
The term $\sum\vec F/m$ represents the acceleration contributed by all external
forces acting on the vehicle. These forces include the aerodynamic force
$\sum\vec F_{\mathrm{aero}}$, the propulsive force
$\sum\vec F_{\mathrm{prop}}$, and the gravitational force
$\sum\vec F_{\mathrm{grav}}$. Their corresponding acceleration vectors are shown
in \cref{fig:force-acceleration-contributions}.

\begin{center}
  \resizebox{0.83\linewidth}{!}{\begin{tikzpicture}[>=Latex,line cap=round,line join=round,scale=0.86]
  % Simplified slender vehicle.
  \draw[line width=0.9pt,rotate=18] (-3.2,-0.32) -- (2.5,-0.32)
    -- (3.0,0) -- (2.5,0.32) -- (-3.2,0.32)
    .. controls (-3.65,0.22) and (-3.65,-0.22) .. cycle;
  \draw[line width=0.8pt,rotate=18] (1.65,0.32) -- (2.15,1.05) -- (2.45,0.38);
  \draw[line width=0.8pt,rotate=18] (1.65,-0.32) -- (2.15,-1.05) -- (2.45,-0.38);
  \coordinate (cg) at (0,0);
  \draw[->,line width=1.05pt] (cg) -- (4.0,1.10)
    node[right,align=left] {$\displaystyle\vec a_{\mathrm{aero}}=
      \frac{\sum\vec F_{\mathrm{aero}}}{m}$};
  \draw[->,line width=1.05pt] (cg) -- (-3.3,-1.55)
    node[left,align=right] {$\displaystyle\vec a_{\mathrm{prop}}=
      \frac{\sum\vec F_{\mathrm{prop}}}{m}$};
  \draw[->,line width=1.05pt] (cg) -- (0,-2.05)
    node[right,align=left] {$\displaystyle\vec a_{\mathrm{grav}}=
      \frac{\sum\vec F_{\mathrm{grav}}}{m}$};
\end{tikzpicture}
}
  {\captionsetup{hypcap=false}\captionof{figure}{Accelerations from forces acting on the vehicle.}}
  \label{fig:force-acceleration-contributions}
\end{center}

Aerodynamic forces are functions of the vehicle geometry, the vehicle state, and
atmospheric properties such as temperature, pressure, and density. The atmosphere
models used to obtain those properties are described in
\cref{sec:atmosphere-model}; aerodynamic-force calculations from input
coefficients are described in \cref{sec:aerodynamic-forces}.

Propulsive forces from rocket or jet engines are defined by input data and are
described in \cref{sec:propulsive-forces}.

TAOS assumes that trajectories remain sufficiently near the earth for terrestrial
gravity to predominate. Gravitational forces from the moon, sun, and planets are
neglected. The earth gravity model used to compute
$\vec a_{\mathrm{grav}}$ is described in \cref{sec:gravity-model}.
